% !TEX root = ../main.tex

\section{Multipole Ring Analysis} \label{sec:multipole_sensitivity}

\subsection{Background of Multipoles Analysis} \label{subsec:BG_multipole}
To address the multi-harmonic characteristic of the propeller noise, a stationary ring of monopoles, referred to as the multipole ring case, is used. This setup replicates the thickness (monopole) noise data first established by Farassat \cite{Farassat1977} for a three bladed propeller, which remains a benchmark for research in aeroacoustics \cite{WoanGregorek1978, Wang2024}. Instead of modelling the blades, discrete point sources form a ring, with a radius equal to the reference propeller. This approach represents the "blade passing points," where the discrete sources represent thickness of the rotor blade tips passing through the ring's azimuthal positions.


\subsection{Mimicking the Noise Signature of a Propeller using Analytical Multipoles} \label{subsec:analytical_multipole}

\subsection{Sensitivity of the FW-H Method in Multipole Ring Prediction}